\section{Moe-эпоха и "database character": персонаж как набор элементов, распознаваемых фанатской культурой}

Здесь сходятся две линии, которые до сих пор шли параллельно. В 1.1 мы видели зрителя, обученного достраивать целое из фиксированных, неполных сигналов. В 1.2 — глаз Тэдзуки, первую деталь лица, которую можно менять независимо от остального и которая одна несёт полное сообщение. Moe-эпоха 2000-х — это момент, когда обе линии становятся не просто творческим приёмом, а рыночным механизмом, и стоит увидеть, как именно это произошло, а не просто зафиксировать факт.

К началу 2000-х у индустрии уже была огромная библиотека наработанных типажей (см. 1.3--1.4), а додзинси-рынок и Comiket дали фанатам инструмент, которого не было у зрителя но или кабуки, — возможность самим рисовать и продавать производные версии персонажей. И вот что обнаружил этот рынок: если тренированный зритель умеет достраивать личность по паре узнаваемых деталей (1.1), а деталь способна нести полное сообщение отдельно от остального лица (1.2), то эти детали в принципе можно отделить от персонажа-источника и продавать поштучно. Кошачьи ушки, meganekko-очки, конкретная реплика в tsundere-регистре, оттенок цвета волос — в moe-жанре вроде "Lucky Star" и чуть позже "K-On!" каждая такая деталь становится единицей, которую фанат может распознать, оценить и захотеть отдельно от остальных, как деталь конструктора, а не как часть кого-то конкретного.

Адзума описывает это как смену модели идентификации: раньше зритель считывал персонажа как целостный grand narrative — историю, из которой характер вырастает. В moe-культуре персонаж читается напрямую через базу распознаваемых элементов, без необходимости в объединяющем нарративе вообще. Отсюда и термин database animals: не "звери" в смысле дикости, а в смысле поведения, отвечающего на стимул-элемент напрямую, без опосредующей личности. Это тот же механизм тренированного считывания, что и в но, доведённый до предела: там зритель достраивал личность из маски и наклона головы, здесь — из ушек и цвета волос, но конвенция сотрудничества та же самая, только личность, которую предстоит достроить, необязательно вообще существует за элементами.

Это не значит, что moe-персонажи обязательно плохи или пусты, — многие отдельные тайтлы этой эпохи прекрасно держат цельный характер. Но культурная механика, которую описал Адзума, реальна и с тех пор никуда не делась: индустрия обучилась производить и продавать узнаваемые элементы отдельно от персонажа, которому они принадлежат. Это ровно тот механизм, о котором предупреждает 0.5 применительно к Рин: собрать её из проверенных элементов легче, чем вырастить из них цельный характер, и рынок вокруг будет систематически подталкивать именно к первому, более лёгкому пути, — потому что зритель для этого пути уже готов и обучен за много веков до появления самого рынка.
